% =====================================================================
%  1.3 Rotuttrykk og brøkeksponenter — fasit
% =====================================================================

\fk{Del 1 — Kvadratrot og $n$-te rot}
\fo{1.1} a) $9$ \ b) $3$ \ c) $-5$ \ d) $\frac{2}{5}$
\fo{1.2} a) $2\sqrt{7}$ \ b) $5\sqrt{2}$ \ c) $3\sqrt{5}$ \ d) $6\sqrt{2}$
\fo{1.3} a) $3\sqrt{2}+5\sqrt{2}-2\sqrt{2} = 6\sqrt{2}$ \ b) $3\sqrt{7}+2\sqrt{7} = 5\sqrt{7}$ \ c) $6\sqrt{5}-3\sqrt{5} = 3\sqrt{5}$ \ d) $\sqrt{36} = 6$
\fo{1.4} a) $2\sqrt{3}+3\sqrt{3} = 5\sqrt{3}$ \ b) $\frac{3\sqrt{3}}{6} = \frac{\sqrt{3}}{2}$ \ c) $\sqrt{2}$ \ d) $\frac{10}{2\sqrt{5}} = \sqrt{5}$
\fo{1.5} Feilen er i første steg: $\sqrt{8}+\sqrt{18} \neq \sqrt{26}$. Riktig er $2\sqrt{2}+3\sqrt{2} = 5\sqrt{2}$, så $\sqrt{2}\cdot 5\sqrt{2} = 10$.
\fo{1.6} a) Noen ganger. For $x \ge 0$ stemmer det. For $x = -3$ er $\sqrt{9} = 3$, ikke $-3$. Riktig er $\sqrt{x^2} = |x|$.
b) Noen ganger. Kvadrering gir $a + b + 2\sqrt{a}\sqrt{b} = a + b$ bare når $\sqrt{a}\sqrt{b} = 0$, altså når $a = 0$ eller $b = 0$.
c) Noen ganger. Venstre side er $2\sqrt{a}$, og $2\sqrt{a} = \sqrt{2a}$ gir ved kvadrering $4a = 2a$, altså $a = 0$.
\fo{1.7} a) $\sqrt{a^2b} = \sqrt{a^2}\cdot\sqrt{b} = a\sqrt{b}$, fordi $a \ge 0$. \ b) $\sqrt{75a^2} = \sqrt{25\cdot3\cdot a^2} = 5a\sqrt{3}$
\fo{Sjekk 1} a) $4\sqrt{2}$ \ b) $5\sqrt{3}-2\sqrt{3} = 3\sqrt{3}$ \ c) Nei. $2\sqrt{3} = \sqrt{4\cdot 3} = \sqrt{12}$, og $\sqrt{12} \neq \sqrt{6}$.

\fk{Del 2 — Brøk som eksponent}
\fo{2.1} a) $7^{\tfrac12}$ \ b) $5^{\tfrac13}$ \ c) $x^{\tfrac32}$ \ d) $x^{\tfrac23}$
\fo{2.2} a) $5$ \ b) $3$ \ c) $8$ \ d) $16$
\fo{2.3} a) $\frac{1}{3}$ \ b) $\frac{1}{27}$ \ c) $\frac{3}{2}$ \ d) $4$
\fo{2.4} a) $x^{\tfrac32}$ \ b) $x^{\tfrac56}$ \ c) $x^{\tfrac74}$ \ d) $x^{\tfrac16}$
\fo{2.5} Feilen er i siste steg: $2^{-2} = \frac{1}{4}$, ikke $-4$. Negativ eksponent gir brøk, ikke negativt tall. Riktig svar: $\frac{1}{4}$.
\fo{2.6} a) $T = 1,52^{\tfrac32} \approx 1,87 \approx 1,9$ år \ b) $a^{\tfrac32} = 8$, så $a = 8^{\tfrac23} = 4$ astronomiske enheter.
\fo{2.7} a) $x = 27^{\tfrac23} = (\sqrt[3]{27})^2 = 9$.
b) Når $n$ er et kvadrattall. Vi har $n^{\tfrac32} = n\sqrt{n}$. Er $n = m^2$, blir det $m^3$, et heltall. Er $n$ ikke et kvadrattall, er $\sqrt{n}$ irrasjonalt, og da er $n\sqrt{n}$ irrasjonalt.
\fo{2.8} a) Aldri. Venstre side er positiv, høyre side er negativ.
b) Noen ganger. $\sqrt{a} = \frac{a}{2}$ gir $a = 4$ (for $a > 0$).
c) Noen ganger. Del med $\sqrt{a} > 0$: $\sqrt{a} < a \iff 1 < \sqrt{a} \iff a > 1$. For $a = 1$ er det likhet. For $a = \frac14$ er $\sqrt{a} = \frac12 > \frac14$.
\fo{Sjekk 2} a) $27$ \ b) $\frac{1}{5}$ \ c) $(a^{\tfrac12})^2 = a^1 = a$. Et positivt tall som ganget med seg selv gir $a$, er $\sqrt{a}$.

\fk{Ekstra utfordring}
\fo{U1} a) $\frac{\sqrt{3}+1}{2}$ \ b) $\sqrt{5}-2$ \ c) $\frac{6(\sqrt{7}+1)}{6} = \sqrt{7}+1$
\fo{U2} a) $7+4\sqrt{3}$ \ b) $5-2 = 3$ \ c) $(5+2\sqrt{6})-(5-2\sqrt{6}) = 4\sqrt{6}$
\fo{U3} a) $2^x = 2^{\tfrac32}$, $x = \frac32$ \ b) $3^{2x} = 3^{\tfrac13}$, $x = \frac16$ \ c) $2^{3x} = 2^{-\tfrac12}$, $x = -\frac16$
\fo{U4} a) Runde 1: $\frac{3}{2}$. Runde 2: $\frac{17}{12}$. b) 4 runder.
c) Er $x$ for stor, er $\frac{2}{x}$ for liten, og $\sqrt{2}$ ligger mellom dem. Gjennomsnittet ligger midt mellom $x$ og $\frac{2}{x}$, og $\sqrt{2}$ ligger mellom de to. Da er gjennomsnittet nærmere $\sqrt{2}$ enn $x$ var.
d) Bytt 2 med 10 i testen og i $2/x$. Det trengs 6 runder.
\fo{U5} Sjette potens fjerner alle rottegnene: $\left(\sqrt{2}\right)^6 = 2^3 = 8$, $\left(\sqrt[3]{3}\right)^6 = 3^2 = 9$ og $\left(\sqrt[6]{6}\right)^6 = 6$. For positive tall beholdes rekkefølgen når vi opphøyer i 6, fordi $x \mapsto x^6$ er økende. Siden $6 < 8 < 9$, er $\sqrt[6]{6} < \sqrt{2} < \sqrt[3]{3}$.
\fo{U6} a) $2$ \ b) Et oddetall i andre er et oddetall, så $p$ må være partall. c) $4k^2 = 2q^2$, så $q^2 = 2k^2$.
d) $q^2$ er partall, så $q$ er partall. Da har $p$ og $q$ felles faktor 2, og det strider mot antakelsen.
e) Steg d) bryter sammen. Vi får $p^2 = 4q^2$ og $p = 2k$, så $q^2 = k^2$. Det sier ikke at $q$ er partall, så selvmotsigelsen uteblir. Det er greit, for $\sqrt{4} = 2 = \frac21$ er rasjonalt.
\fo{U7} Flere svar er riktige. Eksempel: $\frac{4}{\sqrt{7}+\sqrt{3}} = \frac{4(\sqrt{7}-\sqrt{3})}{7-3} = \sqrt{7}-\sqrt{3}$.
Generelt: $\frac{c}{\sqrt{a}+\sqrt{b}} = \frac{c(\sqrt{a}-\sqrt{b})}{a-b}$. Betingelse: $a \neq b$, og $a-b$ må gå opp i $c$.
